\section*{第 \ref{ch:g11:cell-cycle-mitosis} 章——细胞周期与有丝分裂}

\begin{solution}{exo:g11:cell-cycle-mitosis:1}
$\mathrm{G_1}$：生长；S：\omterm{def:g10:universal-dna:information}{DNA} 复制；$\mathrm{G_2}$：
为分裂做准备；M：\omterm{def:g11:cell-cycle-mitosis:cycle}{有丝分裂}（\omterm{def:g10:cells-common-unit:organelle}{细胞核}的分裂）和
胞质分裂（\omterm{def:g10:cells-common-unit:cell}{细胞质}的分裂）。
\end{solution}

\begin{solution}{exo:g11:cell-cycle-mitosis:2}
\omterm{def:g11:cell-cycle-mitosis:chromosome}{染色单体}是复制后\omterm{def:g11:cell-cycle-mitosis:chromosome}{染色体}的两条相同拷贝之一，
\omterm{def:g11:cell-cycle-mitosis:chromosome}{着丝粒}是两者相连的区域。自 S 期结束起
直至后期，一条\omterm{def:g11:cell-cycle-mitosis:chromosome}{染色体}都具有
两条\omterm{def:g11:cell-cycle-mitosis:chromosome}{染色单体}。
\end{solution}

\begin{solution}{exo:g11:cell-cycle-mitosis:3}
前期：\omterm{def:g11:cell-cycle-mitosis:chromosome}{染色体}凝缩，核膜解体。
中期：\omterm{def:g11:cell-cycle-mitosis:chromosome}{染色体}排列在赤道板上。后期：\omterm{def:g11:cell-cycle-mitosis:chromosome}{染色单体}
分开并移向两极。末期：核膜重新形成，
随后发生胞质分裂。
\end{solution}

\begin{solution}{exo:g11:cell-cycle-mitosis:4}
在 S 期，复制进行到一半时。
\end{solution}

\begin{solution}{exo:g11:cell-cycle-mitosis:5}
46 条\omterm{def:g11:cell-cycle-mitosis:chromosome}{染色体}，92 条\omterm{def:g11:cell-cycle-mitosis:chromosome}{染色单体}。
\end{solution}

\begin{solution}{exo:g11:cell-cycle-mitosis:6}
$\mathrm{G_1}$ \qty{11}{h}，S \qty{8}{h}，$\mathrm{G_2}$ \qty{4}{h}，
M \qty{1}{h}。\omterm{def:g10:universal-dna:information}{DNA} 含量为 $2q$ 的\omterm{def:g10:cells-common-unit:cell}{细胞}处于 $\mathrm{G_2}$ 期，或处于
后期结束之前的\omterm{def:g11:cell-cycle-mitosis:cycle}{有丝分裂}期。
\end{solution}

\begin{solution}{exo:g11:cell-cycle-mitosis:7}
中期：40 条\omterm{def:g11:cell-cycle-mitosis:chromosome}{染色体}，80 条\omterm{def:g11:cell-cycle-mitosis:chromosome}{染色单体}。后期：80 条\omterm{def:g11:cell-cycle-mitosis:chromosome}{染色体}
（各含单条\omterm{def:g11:cell-cycle-mitosis:chromosome}{染色单体}），每极 40 条。子\omterm{def:g10:cells-common-unit:cell}{细胞}：40 条\omterm{def:g11:cell-cycle-mitosis:chromosome}{染色体}，40
条\omterm{def:g11:cell-cycle-mitosis:chromosome}{染色单体}。
\end{solution}

\begin{solution}{exo:g11:cell-cycle-mitosis:8}
\omterm{def:g11:cell-cycle-mitosis:cycle}{有丝分裂}指数 $80/800 = 10\%$；\omterm{def:g11:cell-cycle-mitosis:cycle}{有丝分裂}历时 \qty{2.4}{h}。前期
$48/800 \times 24 = \qty{1.44}{h}$；中期 \qty{0.48}{h}（29 分钟）；
后期 \qty{0.18}{h}（11 分钟）；末期 \qty{0.3}{h}（18 分钟）。
\end{solution}

\begin{solution}{exo:g11:cell-cycle-mitosis:9}
没有纺锤体，\omterm{def:g11:cell-cycle-mitosis:chromosome}{染色单体}便无法被拉开：\omterm{def:g10:cells-common-unit:cell}{细胞}停滞在中期，
\omterm{def:g11:cell-cycle-mitosis:chromosome}{染色体}高度凝缩、清晰可见，每条含两条\omterm{def:g11:cell-cycle-mitosis:chromosome}{染色单体}。
正是在这种状态下，\omterm{def:g11:cell-cycle-mitosis:chromosome}{染色体}才适于拍摄，
用于制作\omterm{def:g11:cell-cycle-mitosis:chromosome}{核型}图。
\end{solution}

\begin{solution}{exo:g11:cell-cycle-mitosis:10}
两米长、缠绕在一起的细线无法分拣；凝缩成
几微米长的棒状后，46 条\omterm{def:g11:cell-cycle-mitosis:chromosome}{染色体}便可排列整齐，并被拉开，
而不致断裂或打结。之后它们会解螺旋，
因为\omterm{def:g10:universal-dna:gene}{基因}只能从松散的形态中读取。
\end{solution}

\begin{solution}{exo:g11:cell-cycle-mitosis:11}
分裂 10 次：$2^{10} = 1024$ 个\omterm{def:g10:cells-common-unit:cell}{细胞}。不对：早期胚胎
在两次分裂之间不生长，其周期几乎全是 S 期和 M 期，
几乎没有 $\mathrm{G_1}$ 期。
\end{solution}

\begin{solution}{exo:g11:cell-cycle-mitosis:12}
一个子\omterm{def:g10:cells-common-unit:cell}{细胞}多出 1 条\omterm{def:g11:cell-cycle-mitosis:chromosome}{染色体}，共 47 条；另一个缺少 1 条，共 45 条；
它们的 \omterm{def:g10:universal-dna:information}{DNA} 量为 $q$ 加减一条\omterm{def:g11:cell-cycle-mitosis:chromosome}{染色体}的量。在体\omterm{def:g10:cells-common-unit:cell}{细胞}中，这种错误
只影响数以万亿计的\omterm{def:g10:cells-common-unit:cell}{细胞}中的一个，该\omterm{def:g10:cells-common-unit:cell}{细胞}通常会死亡或被清除；
但若该\omterm{def:g10:cells-common-unit:cell}{细胞}存活并分裂，且多出或缺失的\omterm{def:g11:cell-cycle-mitosis:chromosome}{染色体}使其生长失控，
肿瘤便可能开始形成。
\end{solution}

\begin{solution}{exo:g11:cell-cycle-mitosis:13}
每天约有 $\num{2e11}$ 次分裂，即任一时刻
有 $\num{2e11}$ 个\omterm{def:g10:cells-common-unit:cell}{细胞}处于循环之中（每个\omterm{def:g10:cells-common-unit:cell}{细胞}每天
产生一个新\omterm{def:g10:cells-common-unit:cell}{细胞}）。阻断\omterm{def:g11:cell-cycle-mitosis:cycle}{有丝分裂}会切断供应；红\omterm{def:g10:cells-common-unit:cell}{细胞}寿命为 120 天，
但有些白\omterm{def:g10:cells-common-unit:cell}{细胞}仅能存活数天，因此一周之内白\omterm{def:g10:cells-common-unit:cell}{细胞}数量便会急剧下降，
随之而来的是感染——这是抗癌药物的典型副作用。
\end{solution}

\begin{solution}{exo:g11:cell-cycle-mitosis:14}
在间期，\omterm{def:g11:cell-cycle-mitosis:chromosome}{染色体}解螺旋而不可见；在后期，\omterm{def:g11:cell-cycle-mitosis:chromosome}{染色体}
处于移动状态，相互混杂。在中期，\omterm{def:g11:cell-cycle-mitosis:chromosome}{染色体}高度凝缩、
彼此分离，仍由两条\omterm{def:g11:cell-cycle-mitosis:chromosome}{染色单体}组成，并位于同一平面上，
可依据其大小和形状加以辨认。
\end{solution}

\begin{solution}{exo:g11:cell-cycle-mitosis:15}
处于 $\mathrm{G_1}$ 期的静止肝\omterm{def:g10:cells-common-unit:cell}{细胞}重新进入\omterm{def:g11:cell-cycle-mitosis:cycle}{细胞周期}，反复经历
S 期、$\mathrm{G_2}$ 期和\omterm{def:g11:cell-cycle-mitosis:cycle}{有丝分裂}，直至肝的质量恢复，然后重新回到静止状态。
第 0 天时，几乎所有\omterm{def:g10:cells-common-unit:cell}{细胞}的 \omterm{def:g10:universal-dna:information}{DNA} 量为 $q$；
第 3 天时，许多\omterm{def:g10:cells-common-unit:cell}{细胞}的 \omterm{def:g10:universal-dna:information}{DNA} 量介于 $q$ 与 $2q$ 之间（处于 S 期），
或为 $2q$（处于 $\mathrm{G_2}$ 期或\omterm{def:g11:cell-cycle-mitosis:cycle}{有丝分裂}期）。
\end{solution}

\begin{solution}{pb:g11:cell-cycle-mitosis:1}
\textbf{1.} $100/1000 = 10\%$。

\textbf{2.} $0.10 \times 20 = \qty{2}{h}$。

\textbf{3.} 前期 $60/1000 \times 20 = \qty{1.2}{h}$（72 分钟）；
中期 $0.36$ h（22 分钟）；后期 $0.16$ h（约 10 分钟）；
末期 $0.28$ h（17 分钟）。

\textbf{4.} 后期。\omterm{def:g11:cell-cycle-mitosis:chromosome}{染色单体}已凝缩并相连；
其移动依靠纤维牵拉几微米，
这些纤维在几分钟内即可缩短。

\textbf{5.} 稀有时期仅见于少数\omterm{def:g10:cells-common-unit:cell}{细胞}，计数过少
则误差大；且只有生长区的\omterm{def:g10:cells-common-unit:cell}{细胞}在循环，
其余部位的指数将为零，会稀释数据。

\textbf{6.} 16 条\omterm{def:g11:cell-cycle-mitosis:chromosome}{染色体}，32 条\omterm{def:g11:cell-cycle-mitosis:chromosome}{染色单体}。

\textbf{7.} 32 条\omterm{def:g11:cell-cycle-mitosis:chromosome}{染色体}在移动，每极 16 条。

\textbf{8.} 前期 $2q$；后期全\omterm{def:g10:cells-common-unit:cell}{细胞} $2q$；末期
每个\omterm{def:g10:cells-common-unit:organelle}{细胞核} $q$。

\textbf{9.} S 期占间期的 $7/19$：约 $900 \times 7/19 \approx
330$ 个\omterm{def:g10:cells-common-unit:cell}{细胞}。

\textbf{10.} 后期：\omterm{def:g11:cell-cycle-mitosis:chromosome}{着丝粒}刚刚分裂，16 条
\omterm{def:g11:cell-cycle-mitosis:chromosome}{染色体}的 32 条\omterm{def:g11:cell-cycle-mitosis:chromosome}{染色单体}现已成为 32 条\omterm{def:g11:cell-cycle-mitosis:chromosome}{染色体}，正走向
两极。

\textbf{11.} $\num{20000}/20 = 1000$，即每小时产生 1000 个新\omterm{def:g10:cells-common-unit:cell}{细胞}。

\textbf{12.} 整个生长区每小时 1000 个\omterm{def:g10:cells-common-unit:cell}{细胞}；沿一条\omterm{def:g10:cells-common-unit:cell}{细胞}列
看，生长区在两个方向上各约百个\omterm{def:g10:cells-common-unit:cell}{细胞}宽，故每列每天
约 $\num{24000}/(100 \times 100) \approx 2$--3 个\omterm{def:g10:cells-common-unit:cell}{细胞}，
每个\omterm{def:g10:cells-common-unit:cell}{细胞} \qty{100}{\micro m}，合计仅零点几毫米。实测
\qty{1}{mm} 说明生长区比假设的窄，或\omterm{def:g10:cells-common-unit:cell}{细胞}
伸长更多；数量级是正确的。

\textbf{13.} 一个起点复制 $2 \times 50 \times 7 \times 3600 =
\num{2.5e6}$ 个碱基对；至少需要 $\num{1.6e10}/\num{2.5e6} \approx 6400$
个复制起点。

\textbf{14.} 每个时期都延长，因此每天完成周期的\omterm{def:g10:cells-common-unit:cell}{细胞}减少；
若各时期同比例减慢，则各时期的比例及
\omterm{def:g11:cell-cycle-mitosis:cycle}{有丝分裂}指数均不变。

\textbf{15.} 已过 S 期的\omterm{def:g10:cells-common-unit:cell}{细胞}将在最初数小时内完成
$\mathrm{G_2}$ 期及\omterm{def:g11:cell-cycle-mitosis:cycle}{有丝分裂}；一天之后已无\omterm{def:g10:cells-common-unit:cell}{细胞}能进入\omterm{def:g11:cell-cycle-mitosis:cycle}{有丝分裂}，故
前期至末期均已消失，仅余间期\omterm{def:g10:cells-common-unit:cell}{细胞}，
停滞于 $\mathrm{G_1}$ 期或 S 期起点。

\textbf{16.} $4/1000 \times 20 = \qty{0.08}{h}$，约 5 分钟，
与 10 分钟相比。并不意外：\omterm{def:g10:cells-common-unit:cell}{细胞}仅 4 或 8 个时，计数
主要受偶然因素支配。

\textbf{17.} $12/2000 \times 20 = \qty{0.12}{h}$，约 7 分钟。

\textbf{18.} 人：$0.04 \times 24 \approx \qty{1}{h}$；大蒜
\qty{2}{h}。数量级相近，植物\omterm{def:g10:cells-common-unit:cell}{细胞}略慢。

\textbf{19.} 肿瘤\omterm{def:g10:cells-common-unit:cell}{细胞}持续循环，而多数正常\omterm{def:g10:cells-common-unit:cell}{细胞}
处于 $\mathrm{G_1}$ 期休息，因此任一时刻
处于\omterm{def:g11:cell-cycle-mitosis:cycle}{有丝分裂}的比例更大。阻断纺锤体的药物因此对肿瘤\omterm{def:g10:cells-common-unit:cell}{细胞}
的作用比对正常\omterm{def:g10:cells-common-unit:cell}{细胞}更频繁——但也作用于分裂迅速的正常组织。

\textbf{20.} 周期 20 小时：$\mathrm{G_1}$ 8，S 7，$\mathrm{G_2}$ 4，
\omterm{def:g11:cell-cycle-mitosis:cycle}{有丝分裂}约 2（前期 72 分钟，中期 22 分钟，后期 7--10 分钟，
末期 17 分钟）。后期最短，此时\omterm{def:g11:cell-cycle-mitosis:chromosome}{着丝粒}分裂，
且每极各得到
这 16 条\omterm{def:g11:cell-cycle-mitosis:chromosome}{染色体}中每条的一条\omterm{def:g11:cell-cycle-mitosis:chromosome}{染色单体}。
\end{solution}
